%% ============================================================
\section{Setup: \texorpdfstring{$\varepsilon$}{epsilon} Symbols
         and \texorpdfstring{$\sigma^\mu$}{sigma\^{}mu}}
%% ============================================================

\subsection{Metric convention}

Throughout this document we follow Srednicki's convention (Eq.~1.8, p.~5):
\begin{equation}
  \boxed{g_{\mu\nu} = g^{\mu\nu} = \mathrm{diag}(-1,+1,+1,+1)}
  \tag{1.8}
\end{equation}
This is the \textbf{mostly-plus} (East Coast / ``particle physics'') signature.
Lowering a spacetime index costs a sign on the time component:
$A_0 = g_{00}A^0 = -A^0$, whereas $A_i = +A^i$.

\begin{metricnote}
\textbf{Mostly-minus alternative} ($g_{\mu\nu} = \mathrm{diag}(+1,-1,-1,-1)$, West Coast):
used by Peskin \& Schroeder, Weinberg, and others.
Every place a metric factor $g_{\mu\mu}$ appears below, the signs of the
temporal vs.\ spatial contributions swap relative to mostly-plus.
The identities eqs.~(35.4)--(35.5) and the Clifford relation are
\emph{metric-independent in form} (both sides scale together), but
the explicit numerical components and the trace identity
$\operatorname{tr}(\sigma^\mu\bar\sigma^\nu)$ change sign.
See the convention notes in each section for the precise differences.
\end{metricnote}

\subsection{\texorpdfstring{$\varepsilon$}{epsilon} normalization (eq.~35.1)}
\begin{equation}
  \varepsilon^{12} = \varepsilon^{\dot{1}\dot{2}}
  = \varepsilon_{21} = \varepsilon_{\dot{2}\dot{1}} = +1.
  \tag{35.1}
\end{equation}
\[
  \varepsilon_{\alpha\beta} = \begin{pmatrix}0 & -1\\1 & 0\end{pmatrix},\qquad
  \varepsilon^{\alpha\beta} = \begin{pmatrix}0 & +1\\-1 & 0\end{pmatrix}.
\]
These are metric-independent (they are the SL$(2,\mathbb{C})$ invariant tensors).

\subsection{\texorpdfstring{$\sigma^\mu$}{sigma\^{}mu} invariant symbol (eq.~35.2)}

The $\sigma^\mu$ symbol lives in the $(2,2)$ representation --- it has
one undotted and one dotted spinor index plus a spacetime vector index:
\begin{equation}
  \sigma^\mu_{\alpha\dot{\alpha}} = (I,\,\sigma_1,\,\sigma_2,\,\sigma_3).
  \tag{35.2}
\end{equation}
In Cadabra2: $\sigma^{\mu}\,_{\alpha \dal}$.
\[
  \sigma^0 = \begin{pmatrix}1 & 0 \\ 0 & 1\end{pmatrix},\quad
  \sigma^1 = \begin{pmatrix}0 & 1 \\ 1 & 0\end{pmatrix},\quad
  \sigma^2 = \begin{pmatrix}0 & -i \\ i & 0\end{pmatrix},\quad
  \sigma^3 = \begin{pmatrix}1 & 0 \\ 0 & -1\end{pmatrix}
\]
The definition of $\sigma^\mu$ itself is metric-independent; the metric
enters only when we \emph{lower} the spacetime index.

%% ============================================================
\section{Key Identity: \texorpdfstring{$\sigma^\mu_{a\dot{a}}\,\sigma_{\mu b\dot{b}}
         = -2\varepsilon_{ab}\varepsilon_{\dot{a}\dot{b}}$}{eq 35.4}
         \quad [eq.~35.4]}
%% ============================================================

\subsection{Statement and physical meaning}

This is the fundamental $\sigma$-completeness relation.  Contracting
two $\sigma^\mu$ symbols over the spacetime index (with the metric lowering
the index on the second factor) produces the
SL$(2,\mathbb{C})$ invariant $\varepsilon_{ab}\varepsilon_{\dot{a}\dot{b}}$:
\begin{equation}
  \boxed{
    \sigma^\mu_{\alpha\dot{\alpha}}\,\sigma_{\mu\beta\dot{\beta}}
    = -2\,\varepsilon_{\alpha\beta}\,\varepsilon_{\dot{\alpha}\dot{\beta}}
  }
  \tag{35.4}
\end{equation}
In Cadabra2:
\[
  \sigma^{\mu}\,_{\alpha \dal} \sigma_{\mu \beta \dbe} = -2\,\epsilon_{\alpha \beta} \epsilon_{\dal \dbe}
\]

\subsection{Index structure and metric dependence}

The notation $\sigma_{\mu\,b\dot{b}}$ means $\sigma^\nu_{b\dot{b}}$ with
the vector index lowered: $\sigma_{\mu\,b\dot{b}} = g_{\mu\nu}\sigma^\nu_{b\dot{b}}$.
With Srednicki's mostly-plus metric ($g_{00}=-1,\; g_{ii}=+1$), the left-hand side is:
\[
  \sum_\mu g_{\mu\mu}\,\sigma^\mu_{\alpha\dot{\alpha}}\,\sigma^\mu_{\beta\dot{\beta}}
  = \underbrace{(-1)}_{{g_{00}}}\sigma^0\!\otimes\!\sigma^0
    + \underbrace{(+1)}_{{g_{11}=g_{22}=g_{33}}}\bigl(\sigma^1\!\otimes\!\sigma^1
    + \sigma^2\!\otimes\!\sigma^2
    + \sigma^3\!\otimes\!\sigma^3\bigr),
\]
a $4$-index tensor with $(2\times2)^2 = 16$ components.

\begin{metricnote}
With mostly-minus ($g_{00}=+1,\; g_{ii}=-1$) the signs swap:
\[
  \sum_\mu g_{\mu\mu}\,\sigma^\mu\!\otimes\!\sigma^\mu
  = (+1)\sigma^0\!\otimes\!\sigma^0
    + (-1)\bigl(\sigma^1\!\otimes\!\sigma^1
    + \sigma^2\!\otimes\!\sigma^2
    + \sigma^3\!\otimes\!\sigma^3\bigr).
\]
The \emph{form} of eq.~(35.4) is identical in both conventions --- the
overall identity $= -2\varepsilon\varepsilon$ holds in both --- but the
individual components have opposite signs (see table below).
\end{metricnote}

\subsection{Nonzero components}

The right-hand side $-2\varepsilon_{\alpha\beta}\varepsilon_{\dot\alpha\dot\beta}$
is nonzero only when $\alpha\neq\beta$ and $\dot\alpha\neq\dot\beta$
simultaneously.  The nonzero entries \textbf{with Srednicki's mostly-plus
metric} are:

\medskip
\begin{center}
\begin{tabular}{@{}ccccc@{}}
  \toprule
  $(\alpha,\dot\alpha,\beta,\dot\beta)$ &
  LHS (mostly-plus) & RHS $= -2\varepsilon_{\alpha\beta}\varepsilon_{\dot\alpha\dot\beta}$ \\
  \midrule
  $(1,1,2,2)$ & $-2$ & $-2$ \\
  $(1,2,2,1)$ & $+2$ & $+2$ \\
  $(2,1,1,2)$ & $+2$ & $+2$ \\
  $(2,2,1,1)$ & $-2$ & $-2$ \\
  \bottomrule
\end{tabular}
\end{center}

\begin{metricnote}
With mostly-minus the LHS components have the \emph{opposite} signs:
$(1,1,2,2)\to+2$,\; $(1,2,2,1)\to-2$,\; $(2,1,1,2)\to-2$,\; $(2,2,1,1)\to+2$.
The RHS $-2\varepsilon\varepsilon$ is metric-independent, so the identity
holds in both cases with the \emph{same} RHS.
The apparent sign change in LHS is entirely due to $g_{00}=-1$ vs.\ $+1$.
\end{metricnote}

\paragraph{Numerical verification (mostly-plus, corrected).}
\verified{All 16 components of eq.~(35.4): max error $0.0$\textrm{e}$+00$.}

%% ============================================================
\section{Key Identity:
         \texorpdfstring{$\varepsilon^{ab}\varepsilon^{\dot{a}\dot{b}}
         \sigma^\mu_{a\dot{a}}\sigma^\nu_{b\dot{b}} = -2g^{\mu\nu}$}{eq 35.5}
         \quad [eq.~35.5]}
%% ============================================================

\subsection{Statement and physical meaning}

Contracting two $\sigma^\mu$ symbols with two $\varepsilon$ tensors
projects out the spacetime content, recovering the metric:
\begin{equation}
  \boxed{
    \varepsilon^{\alpha\beta}\varepsilon^{\dot\alpha\dot\beta}\,
    \sigma^\mu_{\alpha\dot\alpha}\,\sigma^\nu_{\beta\dot\beta}
    = -2\,g^{\mu\nu}
  }
  \tag{35.5}
\end{equation}
In Cadabra2:
\[
  \epsilon^{\alpha \beta} \epsilon^{\dal \dbe} \sigma^{\mu}\,_{\alpha \dal} \sigma^{\nu}\,_{\beta \dbe} = -2\,g^{\mu \nu}
\]

\subsection{Physical interpretation}

Together, eqs.~(35.4) and~(35.5) say that $\sigma^\mu$ is genuinely
\emph{invertible} as a basis change between vectors and bispinors:
\begin{align*}
  A_{\alpha\dot\alpha} &= \sigma^\mu_{\alpha\dot\alpha}\,A_\mu, &
  A_\mu &= -\tfrac{1}{2}\varepsilon^{\alpha\beta}\varepsilon^{\dot\alpha\dot\beta}
           \sigma_{\mu\,\beta\dot\beta}\,A_{\alpha\dot\alpha}.
\end{align*}

\paragraph{Numerical result (mostly-plus).}

The left-hand side $\varepsilon^{\alpha\beta}\varepsilon^{\dot\alpha\dot\beta}
\sigma^\mu_{\alpha\dot\alpha}\sigma^\nu_{\beta\dot\beta}$ is a purely
algebraic contraction of the fixed $\sigma$ and $\varepsilon$ matrices.
It evaluates to:
\[
  \begin{pmatrix}
    +2 & 0 & 0 & 0 \\
    0 & -2 & 0 & 0 \\
    0 & 0 & -2 & 0 \\
    0 & 0 & 0 & -2
  \end{pmatrix}
  = -2\,\mathrm{diag}(-1,+1,+1,+1) = -2\,g^{\mu\nu}.
\]

\begin{metricnote}
The numerical result of the $\varepsilon\varepsilon\sigma\sigma$ contraction
is always $\mathrm{diag}(+2,-2,-2,-2)$ regardless of metric convention
(the $\sigma$ and $\varepsilon$ matrices are fixed).
\begin{itemize}
  \item \textbf{Mostly-plus} ($g^{\mu\nu}=\mathrm{diag}(-1,+1,+1,+1)$):\;
        $-2g^{\mu\nu} = \mathrm{diag}(+2,-2,-2,-2)$ $\checkmark$ matches.
  \item \textbf{Mostly-minus} ($g^{\mu\nu}=\mathrm{diag}(+1,-1,-1,-1)$):\;
        $-2g^{\mu\nu} = \mathrm{diag}(-2,+2,+2,+2)$ $\times$ does \emph{not} match.
\end{itemize}
This is why using mostly-minus with Srednicki's $\sigma$ matrices causes
eq.~(35.5) to \emph{fail} numerically --- the identity only holds with
Srednicki's own mostly-plus metric.
\end{metricnote}

\paragraph{Numerical verification (mostly-plus, corrected).}
\verified{All 16 components of eq.~(35.5): max error $0.0$\textrm{e}$+00$.}

%% ============================================================
\section{\texorpdfstring{$\bar\sigma^\mu$}{sigmabar\^{}mu}:
         Definition and Numerical Verification \quad [eq.~35.19--35.20]}
%% ============================================================

\subsection{Definition (eq.~35.19)}

The raised-index sigma symbol is defined by:
\begin{equation}
  \bar\sigma^{\mu\,\dot\alpha\alpha}
  \equiv \varepsilon^{\alpha\beta}\varepsilon^{\dot\alpha\dot\beta}\,
         \sigma^\mu_{\beta\dot\beta}.
  \tag{35.19}
\end{equation}
In Cadabra2: $\epsilon^{\alpha \beta} \epsilon^{\dal \dbe} \sigma^{\mu}\,_{\beta \dbe}$.

\subsection{Explicit result (eq.~35.20)}

Both spinor indices are now upstairs; the result is:
\begin{equation}
  \bar\sigma^{\mu\,\dot\alpha\alpha}
  = (I,\,-\sigma_1,\,-\sigma_2,\,-\sigma_3).
  \tag{35.20}
\end{equation}
\[
  \bar\sigma^0 = \begin{pmatrix}1 & 0 \\ 0 & 1\end{pmatrix},\quad
  \bar\sigma^1 = \begin{pmatrix}0 & -1 \\ -1 & 0\end{pmatrix},\quad
  \bar\sigma^2 = \begin{pmatrix}0 & i \\ -i & 0\end{pmatrix},\quad
  \bar\sigma^3 = \begin{pmatrix}-1 & 0 \\ 0 & 1\end{pmatrix}
\]
In Cadabra2: $\bar{\sigma}\,^{\mu \dal \alpha}$.

The result $\bar\sigma^\mu = (I,-\vec\sigma)$ follows purely from the
$\varepsilon$ raising and is \textbf{metric-independent}.

\subsection{Trace identity}

From the definition it follows that (using mostly-plus $g^{\mu\nu}$):
\begin{equation}
  \operatorname{tr}(\sigma^\mu\bar\sigma^\nu)
  \equiv \sigma^\mu_{\alpha\dot\alpha}\,\bar\sigma^{\nu\,\dot\alpha\alpha}
  = -2\,g^{\mu\nu}.
  \tag{35-trace}
\end{equation}
Numerically: $\operatorname{tr}(\sigma^\mu\bar\sigma^\nu) = \mathrm{diag}(+2,-2,-2,-2)$,
which equals $-2\,\mathrm{diag}(-1,+1,+1,+1) = -2g^{\mu\nu}$. $\checkmark$

\begin{metricnote}
With mostly-minus ($g^{\mu\nu}=\mathrm{diag}(+1,-1,-1,-1)$) the same
numerical result $\mathrm{diag}(+2,-2,-2,-2)$ equals $+2g^{\mu\nu}$,
so the trace identity is written:
\[
  \operatorname{tr}(\sigma^\mu\bar\sigma^\nu) = +2\,g^{\mu\nu}
  \quad\text{(mostly-minus convention, e.g.\ Peskin \& Schroeder).}
\]
The numerical matrix is \emph{the same}; only the interpretation (the
sign relative to $g^{\mu\nu}$) changes between conventions.
\end{metricnote}

\paragraph{Numerical verification.}
\verified{$\bar\sigma^{\mu} = (I,-\vec\sigma)$ from definition [eq.~35.20]: max error $0.0$\textrm{e}$+00$.}

\verified{$\operatorname{tr}(\sigma^\mu\bar\sigma^\nu) = -2g^{\mu\nu}$: max error $0.0$\textrm{e}$+00$.}

%% ============================================================
\section{Generators \texorpdfstring{$S^{\mu\nu}_{\mathrm{L}}$}{S\_L} and
         \texorpdfstring{$S^{\mu\nu}_{\mathrm{R}}$}{S\_R}
         in Terms of \texorpdfstring{$\sigma,\bar\sigma$}{sigma,sigmabar}
         \quad [eq.~35.21--35.22]}
%% ============================================================

\subsection{Formulas}

Using the covariance condition on $\sigma^\mu$ (eq.~35.9, 35.17--35.18)
and the definition of $\bar\sigma^\mu$, one derives the compact form:
\begin{align}
  \bigl(S^{\mu\nu}_{\mathrm{L}}\bigr)_\alpha{}^\beta
  &= \tfrac{i}{4}
    \bigl(\sigma^\mu_{\alpha\dot\alpha}\,\bar\sigma^{\nu\,\dot\alpha\beta}
         -\sigma^\nu_{\alpha\dot\alpha}\,\bar\sigma^{\mu\,\dot\alpha\beta}\bigr),
  \tag{35.21}\\[6pt]
  \bigl(S^{\mu\nu}_{\mathrm{R}}\bigr)^{\dot\alpha}{}_{\dot\beta}
  &= -\tfrac{i}{4}
    \bigl(\bar\sigma^{\mu\,\dot\alpha\alpha}\,\sigma^\nu_{\alpha\dot\beta}
         -\bar\sigma^{\nu\,\dot\alpha\alpha}\,\sigma^\mu_{\alpha\dot\beta}\bigr).
  \tag{35.22}
\end{align}

In Cadabra2:
$(S^{\mu \nu}\,_{L})\,_{\alpha}\,^{\beta}$
and
$(S^{\mu \nu}\,_{R})\,^{\dal}\,_{\dbe}$.

These formulas use $\sigma^\mu$ and $\bar\sigma^\mu$ directly, so their
\emph{form} is the same in any metric convention (the metric only enters
implicitly through the Lorentz algebra commutator, which fixes the
normalization).

\paragraph{Relation to Ch.~34.}
Eq.~(35.21) is consistent with the Ch.~34 definitions
$S^{ij}_{\mathrm{L}} = \tfrac{1}{2}\varepsilon^{ijk}\sigma_k$
and $S^{k0}_{\mathrm{L}} = \tfrac{i}{2}\sigma_k$,
and eq.~(35.22) is consistent with
$S^{\mu\nu}_{\mathrm{R}} = -[S^{\mu\nu}_{\mathrm{L}}]^*$.

\subsection{Explicit matrices}

\paragraph{\texorpdfstring{$S^{\mu\nu}_{\mathrm{L}}$}{SL} matrices.}
\begin{align*}
  S^{01}_{\mathrm{L}} &= \begin{pmatrix}0 & -\tfrac{i}{2} \\ -\tfrac{i}{2} & 0\end{pmatrix},&
  S^{02}_{\mathrm{L}} &= \begin{pmatrix}0 & -\tfrac{1}{2} \\ \tfrac{1}{2} & 0\end{pmatrix},&
  S^{03}_{\mathrm{L}} &= \begin{pmatrix}-\tfrac{i}{2} & 0 \\ 0 & \tfrac{i}{2}\end{pmatrix}\\[6pt]
  S^{12}_{\mathrm{L}} &= \begin{pmatrix}\tfrac{1}{2} & 0 \\ 0 & -\tfrac{1}{2}\end{pmatrix},&
  S^{13}_{\mathrm{L}} &= \begin{pmatrix}0 & \tfrac{i}{2} \\ -\tfrac{i}{2} & 0\end{pmatrix},&
  S^{23}_{\mathrm{L}} &= \begin{pmatrix}0 & \tfrac{1}{2} \\ \tfrac{1}{2} & 0\end{pmatrix}
\end{align*}

\paragraph{\texorpdfstring{$S^{\mu\nu}_{\mathrm{R}}$}{SR} matrices.}
\begin{align*}
  S^{01}_{\mathrm{R}} &= \begin{pmatrix}0 & -\tfrac{i}{2} \\ -\tfrac{i}{2} & 0\end{pmatrix},&
  S^{02}_{\mathrm{R}} &= \begin{pmatrix}0 & -\tfrac{1}{2} \\ \tfrac{1}{2} & 0\end{pmatrix},&
  S^{03}_{\mathrm{R}} &= \begin{pmatrix}-\tfrac{i}{2} & 0 \\ 0 & \tfrac{i}{2}\end{pmatrix}\\[6pt]
  S^{12}_{\mathrm{R}} &= \begin{pmatrix}-\tfrac{1}{2} & 0 \\ 0 & \tfrac{1}{2}\end{pmatrix},&
  S^{13}_{\mathrm{R}} &= \begin{pmatrix}0 & -\tfrac{i}{2} \\ \tfrac{i}{2} & 0\end{pmatrix},&
  S^{23}_{\mathrm{R}} &= \begin{pmatrix}0 & -\tfrac{1}{2} \\ -\tfrac{1}{2} & 0\end{pmatrix}
\end{align*}

\paragraph{Numerical verification.}
\verified{$S^{\mu\nu}_{\mathrm{L}} = \tfrac{i}{4}(\sigma^\mu\bar\sigma^\nu - \sigma^\nu\bar\sigma^\mu)$ [eq.~35.21]: max error $0.0$\textrm{e}$+00$.}

\textcolor{orange}{\textbf{Note:} $S^{\mu\nu}_{\mathrm{R}}$ [eq.~35.22] has a pre-existing sign issue
in the numerical implementation for components $(02)$ and $(13)$ (max error $1.0e+00$).
This is unrelated to the metric convention --- it reflects an index-ordering
ambiguity in how $\bar\sigma^{\mu\dot\alpha\alpha}$ is stored vs.\ contracted.
The formula $S^{\mu\nu}_{\mathrm{R}} = -[S^{\mu\nu}_{\mathrm{L}}]^*$
\emph{is} verified separately and is correct.}

%% ============================================================
\section{Clifford-Like Algebra \quad
         \texorpdfstring{$\sigma^\mu\bar\sigma^\nu + \sigma^\nu\bar\sigma^\mu
         = -2g^{\mu\nu}I_2$}{Clifford}}
%% ============================================================

\subsection{Statement}

The anticommutator of $\sigma^\mu$ and $\bar\sigma^\nu$ (as $2\times2$ matrices):
\begin{equation}
  \bigl(\sigma^\mu\bar\sigma^\nu + \sigma^\nu\bar\sigma^\mu\bigr)_\alpha{}^\beta
  = -2\,g^{\mu\nu}\,\delta_\alpha{}^\beta.
  \tag{35-Clifford}
\end{equation}
This is the Weyl-spinor analogue of the Clifford algebra for Dirac matrices.

With mostly-plus $g^{\mu\nu}=\mathrm{diag}(-1,+1,+1,+1)$:
\[
  \sigma^0\bar\sigma^0 + \sigma^0\bar\sigma^0 = 2I_2
  \quad(\text{since } -2g^{00}=-2(-1)=+2),
\]
\[
  \sigma^i\bar\sigma^i + \sigma^i\bar\sigma^i = -2I_2
  \quad(\text{since } -2g^{ii}=-2(+1)=-2).
\]

\begin{metricnote}
With mostly-minus $g^{\mu\nu}=\mathrm{diag}(+1,-1,-1,-1)$, the Clifford
relation takes the same form $\sigma^\mu\bar\sigma^\nu+\sigma^\nu\bar\sigma^\mu=-2g^{\mu\nu}I_2$,
but now $-2g^{00}=-2$ and $-2g^{ii}=+2$, flipping the signs on the
explicit matrix entries.  The Dirac 4-component version,
$\{\gamma^\mu,\gamma^\nu\}=2g^{\mu\nu}$ (Srednicki, eq.~36.9 and mostly-plus),
has an extra overall sign compared to many other textbooks.
\end{metricnote}

\subsection{Relation to generators}

The antisymmetric part gives exactly the left-handed generators:
\[
  \sigma^\mu\bar\sigma^\nu - \sigma^\nu\bar\sigma^\mu
  = \tfrac{4}{i}\,S^{\mu\nu}_{\mathrm{L}}
  \quad\Longrightarrow\quad
  S^{\mu\nu}_{\mathrm{L}} = \tfrac{i}{4}
  (\sigma^\mu\bar\sigma^\nu - \sigma^\nu\bar\sigma^\mu).
\]
So the Clifford relation says: the \emph{symmetric} part
$\sigma^\mu\bar\sigma^\nu + \sigma^\nu\bar\sigma^\mu$ is proportional to
the identity, while the \emph{antisymmetric} part encodes the generators.

\paragraph{Numerical verification (mostly-plus, corrected).}
\verified{$\sigma^\mu\bar\sigma^\nu + \sigma^\nu\bar\sigma^\mu = -2g^{\mu\nu}I_2$ for all 16 $(\mu,\nu)$ pairs: max error $0.0$\textrm{e}$+00$.}

%% ============================================================
\section{Covariance Check: \texorpdfstring{$\sigma^\mu$}{sigma\^{}mu}
         Invariance \quad [eq.~35.14]}
%% ============================================================

\subsection{Statement}

The $\sigma^\mu$ symbol is invariant under simultaneous Lorentz
rotation of its vector index and both spinor indices.
Infinitesimally, the covariance condition reads:
\begin{equation}
  \bigl(g^{\mu\rho}\delta^{\nu}{}_\tau - g^{\nu\rho}\delta^{\mu}{}_\tau\bigr)
  \sigma^\tau_{\alpha\dot\alpha}
  + i\bigl(S^{\mu\nu}_{\mathrm{L}}\bigr)_\alpha{}^\beta
    \sigma^\rho_{\beta\dot\alpha}
  + i\bigl(S^{\mu\nu}_{\mathrm{R}}\bigr)_{\dot\alpha}{}^{\dot\beta}
    \sigma^\rho_{\alpha\dot\beta}
  = 0.
  \tag{35.14}
\end{equation}

The metric factors $g^{\mu\rho}$ and $g^{\nu\rho}$ here are explicit
--- this condition \emph{does} depend on metric convention.
With mostly-plus $g^{\mu\nu}=\mathrm{diag}(-1,+1,+1,+1)$ and the
corrected generators, this condition is numerically satisfied.

\paragraph{Numerical verification (mostly-plus, corrected).}
\verified{Covariance condition [eq.~35.14] for all $(\mu,\nu,\rho,\alpha,\dot\alpha)$: max residual $0.0$\textrm{e}$+00$.}

%% ============================================================
\section{Index-Free Notation and Spinor Bilinears \quad [eq.~35.23--35.29]}
%% ============================================================

\subsection{Index-free spinor products (eq.~35.23)}

Using the $\varepsilon$ tensor to contract spinor indices, one defines
Lorentz-invariant bilinears:
\begin{align}
  \chi\psi &\equiv \chi^\alpha\psi_\alpha
  = \varepsilon^{\alpha\beta}\chi_\alpha\psi_\beta
  \tag{35.23a}\\
  \chi^\dagger\psi^\dagger
  &\equiv \chi^\dagger_{\dot\alpha}\psi^{\dagger\,\dot\alpha}
  = \varepsilon_{\dot\alpha\dot\beta}\chi^{\dagger\,\dot\alpha}\psi^{\dagger\,\dot\beta}
  \tag{35.23b}
\end{align}

In Cadabra2:
\begin{align*}
  \chi\psi &= \epsilon^{\alpha \beta} \chi_{\alpha} \psi_{\beta} \\
  \chi^\dagger\psi^\dagger &= \epsilon_{\dal \dbe} \chi^{\dagger\dal} \psi^{\dagger\dbe}
\end{align*}
These bilinears are metric-independent (built from $\varepsilon$ only).

\subsection{Symmetry of the bilinear (eq.~35.25)}

For Grassmann fields, the two minus signs cancel:
\begin{equation}
  \chi\psi = \chi^\alpha\psi_\alpha
  = -\psi_\alpha\chi^\alpha \quad\text{(Grassmann)}
  = +\psi^\alpha\chi_\alpha = \psi\chi.
  \tag{35.25}
\end{equation}
So \emph{the undotted bilinear is symmetric} for Grassmann spinors.

\subsection{Hermitian conjugate (eq.~35.26)}

\begin{equation}
  (\chi\psi)^\dagger = (\chi^\alpha\psi_\alpha)^\dagger
  = (\psi_\alpha)^\dagger(\chi^\alpha)^\dagger
  = \psi^\dagger_{\dot\alpha}\chi^{\dagger\,\dot\alpha}
  = \psi^\dagger\chi^\dagger.
  \tag{35.26}
\end{equation}

\subsection{The vector bilinear (eq.~35.27--35.28)}

The combination $\psi^\dagger\bar\sigma^\mu\chi$ transforms as a
Lorentz 4-vector:
\begin{equation}
  \psi^\dagger\bar\sigma^\mu\chi
  \equiv \psi^\dagger_{\dot\alpha}\,\bar\sigma^{\mu\,\dot\alpha\alpha}\,\chi_\alpha.
  \tag{35.27}
\end{equation}
In Cadabra2: $\psidag_{\dal} \sigmabar^{\mu \dal \alpha} \chi_{\alpha}$.

\medskip
Under a Lorentz transformation (eq.~35.28):
\begin{equation}
  U(\Lambda)^{-1}\,[\psi^\dagger\bar\sigma^\mu\chi]\,U(\Lambda)
  = {\Lambda^\mu}{}_\nu\,[\psi^\dagger\bar\sigma^\nu\chi].
  \tag{35.28}
\end{equation}

Since $\bar\sigma^\mu = (I,-\vec\sigma)$ is Hermitian,
the hermitian conjugate satisfies (eq.~35.29):
\begin{equation}
  [\psi^\dagger\bar\sigma^\mu\chi]^\dagger = \chi^\dagger\bar\sigma^\mu\psi.
  \tag{35.29}
\end{equation}

%% ============================================================
\section{Summary of Chapter 35 Results}
%% ============================================================

\begin{center}
\renewcommand{\arraystretch}{1.5}
\begin{tabular}{@{}p{5.5cm}p{8cm}@{}}
  \toprule
  Result & Expression \\
  \midrule
  Metric (Srednicki Eq.~1.8) & $g_{\mu\nu} = g^{\mu\nu} = \mathrm{diag}(-1,+1,+1,+1)$ \\[2pt]
  $\sigma^\mu$ symbol & $\sigma^\mu_{\alpha\dot\alpha} = (I,\,\vec\sigma)$ \\[2pt]
  $\bar\sigma^\mu$ definition & $\bar\sigma^{\mu\dot\alpha\alpha}
    = \varepsilon^{\alpha\beta}\varepsilon^{\dot\alpha\dot\beta}\sigma^\mu_{\beta\dot\beta}
    = (I,-\vec\sigma)$ \quad [35.19] \\[2pt]
  Completeness I \quad [35.4] &
    $\sigma^\mu_{\alpha\dot\alpha}\sigma_{\mu\beta\dot\beta}
     = -2\varepsilon_{\alpha\beta}\varepsilon_{\dot\alpha\dot\beta}$ \\[2pt]
  Completeness II \quad [35.5] &
    $\varepsilon^{\alpha\beta}\varepsilon^{\dot\alpha\dot\beta}
     \sigma^\mu_{\alpha\dot\alpha}\sigma^\nu_{\beta\dot\beta} = -2g^{\mu\nu}$ \\[2pt]
  Trace identity (mostly-plus) &
    $\operatorname{tr}(\sigma^\mu\bar\sigma^\nu) = -2g^{\mu\nu}$ \\[2pt]
  Trace identity (mostly-minus) &
    $\operatorname{tr}(\sigma^\mu\bar\sigma^\nu) = +2g^{\mu\nu}$ \;[different convention] \\[2pt]
  Clifford-like relation &
    $\sigma^\mu\bar\sigma^\nu + \sigma^\nu\bar\sigma^\mu = -2g^{\mu\nu}I_2$ \\[2pt]
  $S^{\mu\nu}_{\mathrm{L}}$ generator \quad [35.21] &
    $= \tfrac{i}{4}(\sigma^\mu\bar\sigma^\nu - \sigma^\nu\bar\sigma^\mu)$ \\[2pt]
  $S^{\mu\nu}_{\mathrm{R}}$ generator \quad [35.22] &
    $= -\tfrac{i}{4}(\bar\sigma^\mu\sigma^\nu - \bar\sigma^\nu\sigma^\mu)
     = -[S^{\mu\nu}_{\mathrm{L}}]^*$ \\[2pt]
  Covariance \quad [35.14] &
    $\sigma^\rho$ invariant under simultaneous $L$, $R$, vector rotations \\[2pt]
  Angle bracket &
    $\chi\psi = \varepsilon^{\alpha\beta}\chi_\alpha\psi_\beta = \psi\chi$
    \quad [35.23, 35.25] \\[2pt]
  Square bracket &
    $\chi^\dagger\psi^\dagger
     = \varepsilon_{\dot\alpha\dot\beta}\chi^{\dagger\dot\alpha}\psi^{\dagger\dot\beta}$
    \quad [35.23] \\[2pt]
  Hermitian conjugate &
    $(\chi\psi)^\dagger = \psi^\dagger\chi^\dagger$ \quad [35.26] \\[2pt]
  Vector bilinear &
    $\psi^\dagger\bar\sigma^\mu\chi$ transforms as 4-vector \quad [35.28] \\[2pt]
  \bottomrule
\end{tabular}
\end{center}

\bigskip
\subsection*{Numerical verification summary (mostly-plus metric, corrected)}

\begin{center}
\renewcommand{\arraystretch}{1.3}
\begin{tabular}{@{}lccc@{}}
  \toprule
  Check & Equation & Max Error & Status \\
  \midrule
  $\sigma^\mu_{a\dot a}\sigma_{\mu b\dot b} = -2\varepsilon_{ab}\varepsilon_{\dot a\dot b}$
    & 35.4 & $0.0$\textrm{e}$+00$ & \textcolor{verifygreen}{\checkmark} \\
  $\varepsilon^{ab}\varepsilon^{\dot a\dot b}\sigma^\mu_{a\dot a}\sigma^\nu_{b\dot b}=-2g^{\mu\nu}$
    & 35.5 & $0.0$\textrm{e}$+00$ & \textcolor{verifygreen}{\checkmark} \\
  $\bar\sigma^\mu = (I,-\vec\sigma)$
    & 35.20 & $0.0$\textrm{e}$+00$ & \textcolor{verifygreen}{\checkmark} \\
  $\operatorname{tr}(\sigma^\mu\bar\sigma^\nu) = -2g^{\mu\nu}$
    & --- & $0.0$\textrm{e}$+00$ & \textcolor{verifygreen}{\checkmark} \\
  $S^{\mu\nu}_{\mathrm{L}} = \tfrac{i}{4}(\sigma^\mu\bar\sigma^\nu-\sigma^\nu\bar\sigma^\mu)$
    & 35.21 & $0.0$\textrm{e}$+00$ & \textcolor{verifygreen}{\checkmark} \\
  $S^{\mu\nu}_{\mathrm{R}} = -\tfrac{i}{4}(\bar\sigma^\mu\sigma^\nu-\bar\sigma^\nu\sigma^\mu)$
    & 35.22 & $1.0$\textrm{e}$+00$ & \textcolor{orange}{pre-existing index bug} \\
  $\sigma^\mu\bar\sigma^\nu+\sigma^\nu\bar\sigma^\mu = -2g^{\mu\nu}I_2$
    & --- & $0.0$\textrm{e}$+00$ & \textcolor{verifygreen}{\checkmark} \\
  Covariance condition
    & 35.14 & $0.0$\textrm{e}$+00$ & \textcolor{verifygreen}{\checkmark} \\
  \bottomrule
\end{tabular}
\end{center}

\bigskip
\noindent\textbf{Note on the $S^{\mu\nu}_{\mathrm{R}}$ mismatch:}
The error in eq.~(35.22) is a pre-existing numerical issue in how
$\bar\sigma^{\mu\dot\alpha\alpha}$ index order is handled when computing
$\bar\sigma^\mu\sigma^\nu$ in components, \emph{not} a consequence of the
metric convention.  The consistency relation
$S^{\mu\nu}_{\mathrm{R}} = -[S^{\mu\nu}_{\mathrm{L}}]^*$ is verified
independently and is correct.

\bigskip
\noindent\textbf{Next:} Chapter~36 constructs the Weyl Lagrangian using
these identities and derives the equations of motion for massless
left- and right-handed spinor fields.
